\paragraph{Independent perturbations away from symmetry.}
The covariance range above preserves an invariant initialization family. The next result relaxes the initialization itself at one fixed, nondegenerate configuration; it does not simultaneously perturb data or augmentation.

\begin{theorem}[A parameter neighborhood with encoder rebound]
\label{thm:asymmetric-rebound}
Fix $\Sigma_x=I_2$, independent view noises $\xi_j\sim N(0,I_2/10)$, $\eta_h=\eta_W=1$, $\epsilon=1/100$, and $w_0=1/10$. In the $e_+,e_-$ coordinates put
\begin{equation}
 \bar Q=\begin{pmatrix}1/10&1/\sqrt{200}\\1/10&-1/\sqrt{200}\end{pmatrix},
 \qquad \bar W=(1/10,1/10).
\end{equation}
Allow all six initialization coordinates to vary independently:
\begin{align}
 \max_{c,j}|q_{cj}(0)-\bar Q_{cj}|&\le8\,10^{-6},\nonumber\\
 \max_c|w_c(0)-1/10|&\le8\,10^{-6}.
 \label{eq:asymmetric-cube}
\end{align}
Let $\alpha_c$ be the continuous filter angle, $d=(\alpha_1-\alpha_2)/2$, and let $t_j$ be the first hit of the fixed section
\begin{equation}
 d_j=\arctan\sqrt{\tan\theta_j},\qquad
 \theta_1=3\theta_0/4,\quad \theta_2=\theta_0/6.
\end{equation}
For every task in Theorem~\ref{thm:bottleneck-family}, both encoder and projection risks satisfy
\begin{align}
 R(t_1)&\le R(0)-0.001675,\\
 R(t_2)&\ge R(0)+0.010905.
\end{align}
The two sections are reached with $0<t_1<t_2<900$. Until $t_2$, the ranks remain two and one and the pretraining loss strictly decreases.
\end{theorem}

\paragraph{Geometry outside the invariant family.}
The sections are common to the entire box; physical hitting times may differ. The proof controls filter-angle asymmetry and head imbalance, rather than demanding that all growing parameters remain close. Write $v_c=\norm{q_c}^2$, $m=(v_1+v_2)/2$, $u=(w_1+w_2)/2$, $a=(\alpha_1+\alpha_2)/2$, and
\begin{equation}
 g=\frac{v_1-v_2}{v_1+v_2},\quad z_w=\frac{w_1-w_2}{w_1+w_2},\quad
 \sigma=\tfrac12\log\frac{\tan(2d(0))}{\tan(2d)}.
\end{equation}
Put $D=2d$. The exact angular equations in an increasing clock $s$ are
\begin{align}
 \frac{dd}{ds}&=-u^2m(1-z_w^2)\sin D\cos D,\nonumber\\
 \frac{da}{ds}&=u^2m(1-z_w^2)g\sin D\cos D.
\end{align}
Thus $da/dd=-g$: length imbalance directly rotates the filter midline. Head imbalance also matters, through the evolution of $g$, progress speed, and projection direction. It is not merely a head-selection perturbation.

The encoder space at fixed $d$ is a rotation by $a$ of the symmetric space, regardless of the two positive filter lengths. Its projector therefore obeys
\begin{equation}
 \norm{P_{\rm enc}(a,d)-P_{\rm enc}(0,d)}_2\le2|a|.
\end{equation}
For the projection an additional mixing coefficient appears:
\begin{equation}
 \delta_m=\frac{g+z_w}{1+gz_w},\qquad
 \norm{P_{\rm proj}-\bar P_{\rm proj}}_2\le\sqrt2|a|+|\delta_m|.
\end{equation}
Here the bar denotes the symmetric reference at the same $d$. These bounds target the task spaces directly. Large common growth of the filter lengths or heads need not produce equally large risk error.

The encoder/head distinction remains exact away from symmetry. Let $c_d=\cos d$, $s_d=\sin d$, $\ell_d=\sqrt{c_d^4+s_d^4}$, and let $R_a$ denote input rotation by $a$. The matrices
\begin{align}
 D_a&=R_a\diag(c_d^2,s_d^2)R_a^\top/\ell_d,\nonumber\\
 X_a&=R_aE_\times R_a^\top,\nonumber\\
 N_a&=R_a\diag(s_d^2,-c_d^2)R_a^\top/\ell_d
\end{align}
form an orthonormal basis of the three-dimensional moment space. For a unit quadratic task, write $b_D=\langle\beta,D_a\rangle$, $b_X=\langle\beta,X_a\rangle$, and $b_N=\langle\beta,N_a\rangle$. With $\chi=\sqrt2\delta_m c_ds_d/\ell_d$, direct orthogonal projection gives
\begin{align}
 R_{\rm enc}&=b_N^2,\nonumber\\
 R_{\rm proj}&=b_N^2+\frac{(b_X-\chi b_D)^2}{1+\chi^2}.
 \label{eq:asymmetric-risk-split}
\end{align}
The first term is the unreadable encoder direction; the nonnegative second term is additional head selection. At symmetry $a=\chi=0$, it reduces to the fixed $z^2$ gap in Equation~\eqref{eq:analytic-risks}. Under perturbation that gap can change, so the stability proof bounds both task projectors instead of assuming the symmetric gap remains fixed.

\paragraph{Progress and the retained margin.}
The progress clock removes the factor $u^2m(1-z_w^2)$, leaving $dd/d\sigma=-\sin D\cos D$ and $da/d\sigma=g\sin D\cos D$. The exact nonlinear imbalance equations imply
\begin{align}
 D^+(|g|+|z_w|)&\le4.002(|g|+|z_w|),\nonumber\\
 \left|\frac{da}{d\sigma}\right|
 &\le\tfrac12(|g|+|z_w|).
\end{align}
Here $D^+$ denotes the upper right derivative with respect to $\sigma$. With $e=8\,10^{-6}$ and $N=|g|+|z_w|$, the source box gives $N(0)<33e$ and $|a(0)|<12e$. Hence
\begin{align}
 N(\sigma)&\le33e\exp(4.002\sigma),\nonumber\\
 |a(\sigma)|&\le12e+\frac{33e}{8.004}\big(e^{4.002\sigma}-1\big).
\end{align}
The first and second sections satisfy $\sigma<0.203$ and $\sigma<0.774$. This bounds transverse amplification in a short directional clock, although physical arrival may take hundreds of time units.

Table~\ref{tab:asymmetric-budget} displays the resulting operator-norm projector errors relative to the symmetric reference. Risk differences must include the source error as well as the endpoint error: the comparison is against each perturbed trajectory's own initialization. For both representations, the risk-increment errors are at most $174e$ and $947e$. Subtracting these from the symmetric gaps leaves improvement $>0.001958$ and rebound $>0.014234$, implying the more conservative theorem constants.
\begin{table}[t]
\centering\small
\caption{Uniform projector-error bounds divided by the initialization radius $e$. Source and endpoint errors both enter the risk-increment budget. The columns use common geometric sections, not common physical times.}
\label{tab:asymmetric-budget}
\begin{tabular}{lrrr}
\toprule
Representation&Source&First section&Second section\\\midrule
Encoder&48&34.4&198.4\\
Projection&75&99&872\\\bottomrule
\end{tabular}
\end{table}

The argument also closes continuation. Negative conserved quantities $v_c-2w_c^2$ control relative growth; bounded imbalances keep the heads positive and the two filters independent. An energy barrier $V<0.151$ retains the descending objective regime, and $d\sigma/dt>0.000869$ proves arrival before 900. These checks exclude exits through rank loss, gain loss, or parameter blow-up before either section (Appendix~\ref{app:asymmetric-stability}).

The box has absolute radius $8\,10^{-6}$, about $0.008$--$0.011\%$ of the nonzero base coordinates. Within this small region, filter and head asymmetries alter the task spaces and hitting times, yet every initialization retains the same positive risk gaps. Exact symmetry enables the closed solution; the projector bounds carry its task consequence beyond that symmetry.
